\section{探査データ：測る、整える、地下へ戻す}
\label{theory:exploration:data}
受振器に入る波形は地下構造の写しそのものではない。震源波形、伝播経路、反射・屈折、計器の応答、必要でない波や環境雑音が重なった記録である。計測配置を設計し、時間と振幅をそろえ、対象成分を取り出して地下モデルへ結び付ける。フィルタ、重合、走時補正、画像化、逆解析はそれぞれ異なる役割を持つ。

\subsection{線形性と時不変性：応答を重ねる条件}
\label{theory:exploration:lti}
線形性は入力の和や定数倍に対する応答が同じ操作を保つことである。時不変性は入力を時間移動したとき、応答も同じだけ移動することである。両方が成立するLTI系では単位インパルス応答 \(h\) を使い、
\[
 y(t)=\int_{-\infty}^{\infty}x(\tau)h(t-\tau)\,d\tau,
 \qquad y[n]=\sum_k x[k]h[n-k]
\]
と表せる\cite{mit_signals2010}。任意入力を小さなインパルスの和へ分解し、その応答を加えるという意味である。

時変の線形系でも重ね合わせはできるが、一般には入力時刻にも依存する \(h(t,\tau)\) が必要である。「線形」だけから単一応答の時間移動による畳み込みまで保証しない。計器の応答を補償するときも、応答の零点付近では雑音を強く増幅する場合がある。

\subsection{フィルタ：残す周波数と抑える周波数}
\label{theory:exploration:filters}
低域通過は低周波数を残し高周波数を抑える。\(T=1/f\) なので、長周期を残す処理である。帯域通過は範囲内を残す。ノッチは狭い帯域を阻止する処理であり、商用電源等の周期的混入に用いる\cite{scipy_notch}。ノッチを「特定周波数だけを通す」と覚えると保持と除去が逆になる。

理想フィルタの通過利得1・阻止利得0はモデルの定義である。有限長の実装には遷移帯域、リップル、位相変化、端の処理がある。抑えたい雑音と対象信号の帯域が重なれば、雑音だけを無条件に取り除くことはできない。二次元の画像には空間方向の平滑化や微分を行い、エッジや構造を調べる。フィルタは二次元だけに限定された概念ではない。

\subsection{移動平均：分散の低減と周期消去の反例}
\label{theory:exploration:moving}
\(M\) 点の等重み平均 \(y[n]=M^{-1}\sum_{k=0}^{M-1}x[n-k]\) は平滑化である。各点に独立・平均ゼロ・分散 \(\sigma^2\) の雑音が含まれると、平均後の雑音分散は \(\sigma^2/M\)、標準偏差は \(\sigma/\sqrt{M}\) となる。有限点でランダム雑音が全て消えるわけではない。中央窓は未来側の値も使うため、即時の処理と後処理を区別する。

周波数応答は等比級数を足して独自に導くと、
\[
 H(e^{i\omega})=\frac1M\sum_{k=0}^{M-1}e^{-i\omega k}
 =e^{-i\omega(M-1)/2}\frac{\sin(M\omega/2)}{M\sin(\omega/2)}.
\]
\(\omega=2\pi j/M\) の非DC成分では零点を持つ。このため、周期がちょうど \(M\) 標本の正弦波は消せる。\(M=4\) なら一周期の四つの値の和がゼロとなる。一方、周期がずれれば同じ窓で完全には消えない。MITの移動平均の応答図も周波数によって利得と零点が変わることを示す\cite{mit_moving2012}。平成28年度問20の公式2を、移動平均では周期雑音を決して除去できないという数学的命題へ変更しない。

\subsection{スタックの平方根則を信号と分散から導く}
\label{theory:exploration:stack}
同じ対象信号 \(s\) と雑音 \(n_j\) の記録を \(N\) 本重ねる。
\[
 \sum_{j=1}^{N}(s+n_j)=Ns+\sum_jn_j,
 \qquad \operatorname{Var}\!\left(\sum_jn_j\right)=N\sigma^2
\]
は無相関・同分散雑音の条件で成り立つ。信号振幅は \(N\) 倍、雑音標準偏差は \(\sqrt N\) 倍なので、振幅S/Nの改善倍率は \(\sqrt N\) である。平均として \(N\) で割っても比は同じである。電力S/Nなら改善倍率は \(N\) となる。\(N^2\) はどちらの改善倍率でもない。

一般の雑音では分散に \(2\sum_{i<j}\operatorname{Cov}(n_i,n_j)\) が加わる。同じ電源雑音が各記録に残れば、回数だけ増やして平方根則を期待できない。対象波形も時刻がずれれば相殺されるため、走時差や極性をそろえる。USGSの実データ処理ではNMOとCMP重合、重合後の処理が別々に記録される\cite{usgs_processing2000}。

受振器群は独立雑音の平均化に加え、入射方向と空間波長によるフィルタとして働く。表面波は空間的に相関した波であり、単純な独立雑音とは異なる。群間隔と配列長によって対象反射も弱まる場合がある。同じ起振点で反復する垂直スタックと、共通中点の異なる経路をそろえるCMP重合を区別する。

\subsection{標本化：二倍という数と境界条件}
\label{theory:exploration:sampling}
標本化は \(x[n]=x(nT_s)\)、標本化周波数は \(f_s=1/T_s\) である。帯域制限信号を理想的に再構成するには、周期的に複写されるスペクトルが重ならないことが必要である。最高周波数成分を含む信号には \(f_s>2f_{\max}\) と余裕を取る\cite{mit_signals2010}。

等号では \(x(t)=\sin(2\pi f_{\max}t)\) を \(t=n/(2f_{\max})\) で測ると全てゼロになる。最高周波数を含まない帯域や境界成分の扱いを定義しないまま、「二倍以上だから全て完全復元」と断定しない。現実には量子化、有限記録長、時刻誤差、帯域外雑音もある。aliasingは異なる連続信号が同じ標本を与える現象であり、取得後にFFTをかけるだけでは区別できない。

\subsection{FFT・相互相関・コヒーレンス}
\label{theory:exploration:fft}
有限標本のDFTは \(X[k]=\sum_{n=0}^{N-1}x[n]e^{-2\pi i kn/N}\) である。FFTは同じ和を構造を利用して速く計算する算法である。radix-2では二のべき乗の標本数が扱いやすいが、それ以外が計算不能ではない\cite{scipy_fft}。ゼロを追加して計算点を増やしても、実測時間が延びたり近い二周波数を区別する情報が増えたりするわけではない。

\label{theory:exploration:correlation}
相互相関 \(R_{xy}[\ell]=\sum_n x[n]y[n+\ell]\) は一方をずらして類似を測る。採用する符号 convention と規格化によって遅れの符号や振幅を確認する。周期信号なら複数のピークがあり、最大値だけで唯一の到達遅れを保証しない。

代表的な二乗コヒーレンスは \(C_{xy}(f)=|S_{xy}(f)|^2/(S_{xx}(f)S_{yy}(f))\) であり、周波数ごとの関係を測る\cite{scipy_coherence}。相互相関と関連するが、時間遅れを横軸にする同じ量ではない。非零の自己スペクトルと適切な統計推定の条件を考え、二信号の因果関係をこの値だけで断定しない。

\subsection{走時・静補正・反射・CMPを結ぶ}
\label{theory:exploration:seismicterms}
単純な水平二層の屈折波では走時直線の切片を原点走時と呼ぶ。観測距離ゼロの実測初動と同一ではなく、直線の外挿である。中間層が薄く、初動として一度も現れない場合がある。また、低速度層は通常の臨界屈折条件を満たさず検出しにくい。USGSはこれらを別々に説明する\cite{usgs_refraction2013}。ブラインド層の日本語の使用範囲が広い場合もあるため、層が見えない仕組みを明示する。

単純な垂直入射の反射では \(Z=\rho V\)、\(R=(Z_2-Z_1)/(Z_2+Z_1)\) が反射振幅を支配する。\(Z\) はインピーダンス、\(R\) は反射係数であり別の量である\cite{epa_seismic_reflection}。静補正は地形や表層の低速度層等が生む走時ずれをそろえる\cite{usgs_static2004}。ブーゲー補正は質量による重力を扱い、走時補正ではない。

CMP重合では共通中点の記録を集め、通常NMO等の時間差を補正する。水平・均質の近似から外れると、共通中点が同じ反射点を意味するとは限らない。速度モデル、走時と深度の対応、データの範囲を確かめる。

\subsection{下方接続とmigrationは地震波にも用いる}
\label{theory:exploration:continuation}
migrationは傾斜面の反射や回折記録を地下の位置へ再配置する画像化である。Stanfordの説明では地表の波動場を深度方向へ下方接続し、適切な時刻の波動場から画像を作る\cite{stanford_downward2001}。重力・磁気の下方接続と名前を共有するが、伝播させる場と方程式が異なる。下方接続をポテンシャル場専用の処理と覚えるのは誤りである。

令和元年度問34の初回5はこの限定から生じた。初回記録は保存し、波動下方接続の一次説明と、送受信コイルを用いるrigid-boom EMIの電磁誘導の原理を別に確認して4へ訂正した\cite{rigid_boom2019}。後者は一次論文の検索収録を確認した範囲であり、PDF全文の取得・読了には至っていない。装置の日本語名の最終照合は査読待ちとする。

\subsection{順問題・逆問題・正則化：何を更新するのか}
\label{theory:exploration:inverse}
順問題はモデル \(m\) を与えてデータ \(F(m)\) を計算する。逆問題はデータ \(d\) からモデルを推定する。計測はデータの取得であり、計測そのものを順問題の数値計算と同一視しない。代表的な目的関数は
\[
 \Phi(m)=\|W_d(F(m)-d)\|_2^2+\lambda\|L(m-m_0)\|_2^2
\]
である。最初の項は誤差で重み付けした残差、第二項は滑らかさや参照モデルからの差への拘束である\cite{pygimli_inversion}。データを誤差より細かく再現しても、雑音をモデルへ取り込んでいる可能性がある。

Marquardt型の減衰付き更新、共役勾配、遺伝的探索はモデル更新の算法である。正則化・平滑化拘束は逆解の安定性を制御する。有限要素法や差分法は支配方程式の離散化を担い、\(F(m)\) や感度を計算するため逆解析の内部でも利用できる。平成30年度問35の公式5は直接の役割の比較として読むが、FEMがインバージョンに無関係という絶対命題にはしない。

\subsection{波動の基本則と等価層・検層}
\label{theory:exploration:principles}
線形波動では重ね合わせが成立する。相反性を持つ媒体では震源と受振点を交換した応答に対応がある。ホイヘンスは波面構成、スネルは境界の位相整合による角度と波速、フェルマは波線の走時の停留条件を表す。最短時間という単語だけで複数経路の可能性を排除しない。レンツの法則は磁束変化を妨げる電磁誘導の向きに関する法則であり、弾性波の屈折則とは異なる。

USGSの等価層は重力の密度層や磁気の双極子層によってポテンシャル場を表す説明である\cite{usgs_equivalent2007}。場の表現を波動伝播の基本則と混同しない。ラテラル法は孔内の電極間隔と配置による電気検層の方式である\cite{usgs_lateral_logging}。計測配置の名称と、取得後のフィルタ・差分・最適化の算法を区別する。
